% Generated draft: requires mathematical and source review.
\documentclass[11pt]{article}
\usepackage[margin=25mm]{geometry}
\usepackage{fontspec}
\setmainfont{Latin Modern Roman}
\usepackage{amsmath,amssymb,mathtools,unicode-math}
\setmathfont{Latin Modern Math}
\usepackage{enumitem,needspace,xurl,hyperref}
\setlength{\parindent}{0pt}
\setlength{\parskip}{5pt}
\setlength{\emergencystretch}{3em}
\newcommand{\N}{\mathbb{N}}
\newcommand{\Z}{\mathbb{Z}}
\newcommand{\Q}{\mathbb{Q}}
\newcommand{\R}{\mathbb{R}}
\newcommand{\C}{\mathbb{C}}
\newcommand{\F}{\mathbb{F}}
\newcommand{\A}{\mathbb{A}}
\newcommand{\PP}{\mathbb{P}}
\newcommand{\E}{\mathbb{E}}
\newcommand{\eps}{\varepsilon}
\newcommand{\dd}{\,\mathrm{d}}
\begin{document}
\section*{UnsolvedMath ID: 700722}
\textbf{Source problem code:} catalogue entry
\subsection{Definitions and mathematical statement}
Let \(F\) be a free group, with no restriction on its rank stated in the source. Write \(\Z F\) for its integral group ring and identify \(F\) with its natural image in the group of units of \(\Z F\). Define the augmentation map and augmentation ideal by
\[
\varepsilon:\Z F\longrightarrow\Z,\qquad
\varepsilon\!\left(\sum_{g\in F}a_g g\right)=\sum_{g\in F}a_g,
\qquad \mathfrak f=\ker\varepsilon.
\]
For a normal subgroup \(T\trianglelefteq F\), let
\[
\mathfrak t=\Z F(T-1)=\left\langle t-1\mid t\in T\right\rangle_{\text{two-sided ideal}}.
\]
In particular, use \(\mathfrak r,\mathfrak r_i,\mathfrak s\) for the ideals corresponding to \(R,R_i,S\), respectively. For an ideal \(I\subseteq\Z F\), put
\[
1+I=\{1+a\mid a\in I\},\qquad
D_F(I)=F\cap(1+I)=\{g\in F\mid g-1\in I\}.
\]
Ideal products are taken in their displayed order, and \(\mathfrak f^0=\Z F\). For subgroups \(A,B\leq F\), let \([A,B]\) be the subgroup generated by \(a^{-1}b^{-1}ab\), with \(a\in A\) and \(b\in B\). Write \(R'=[R,R]\), and define the lower central series by \(\gamma_1(F)=F\) and \(\gamma_{j+1}(F)=[\gamma_j(F),F]\).

The source gives the illustrative identification
\[
D_F(\mathfrak r\mathfrak f)=R'.
\]
Its remaining identification requests are the following, uniformly over the free groups and normal subgroups indicated. The source does not explicitly specify the range of the exponent \(n\) in parts (c), (d), and (f); positive integral exponents are the conventional reading, but the intended lower bounds remain a source qualification.
\begin{enumerate}
\item[(a)] For a positive integer \(n\) and normal subgroups \(R_1,\ldots,R_n\trianglelefteq F\), identify, in terms of the corresponding subgroups of \(F\),
\[
D_F(\mathfrak r_1\mathfrak r_2\cdots\mathfrak r_n).
\]
\item[(b)] For normal subgroups \(R_1,R_2,R_3\trianglelefteq F\), identify
\[
D_F(\mathfrak r_1\mathfrak r_2\mathfrak r_3).
\]
The source lists this case separately, although it is contained in (a).
\item[(c)] For \(S\trianglelefteq F\) such that \(F/S\) is finitely generated and nilpotent, identify
\[
D_F(\mathfrak f\mathfrak s+\mathfrak f^n).
\]
\item[(d)] Using the ideal \(\mathfrak s\) corresponding to \(S\trianglelefteq F\), identify
\[
D_F(\mathfrak f\mathfrak s\mathfrak f+\mathfrak f^n).
\]
The source does not explicitly say whether the hypothesis that \(F/S\) is finitely generated and nilpotent carries over from (c). Thus the intended scope of this request, either for all normal \(S\) or only for those satisfying that hypothesis, remains unresolved.
\item[(e)] For \(R\trianglelefteq F\) and integers \(n>k\geq2\), define
\[
\mathfrak r(k)=\sum_{j=0}^{k-1}\mathfrak f^j\mathfrak r\mathfrak f^{k-1-j}.
\]
Identify
\[
D_F\bigl(\mathfrak r(k)+\mathfrak f^n\bigr).
\]
\item[(f)] The source separately archives the following question as solved: for \(R\trianglelefteq F\), is
\[
\frac{D_F(\mathfrak r+\mathfrak f^n)}{R\,\gamma_n(F)}
\]
always an abelian group? Here the denominator is the product subgroup \(R\gamma_n(F)\). The supplied archive reports an affirmative answer; this is a supplied historical assessment, not an independently verified result in this draft.
\end{enumerate}
The requested identifications in (a)--(e) are subgroup descriptions analogous to the displayed example. The supplied record classifies the combined problem as partially solved, with (f) archived as solved and the remaining parts retained in the main section.
\subsection{Short English statement}
For a free group, describe the subgroups obtained by intersecting it with one plus specified products or sums of ideals in its integral group ring. The source also includes an archived question about whether a related quotient is always abelian, reported there as answered affirmatively. The exponent ranges and the scope of the nilpotence hypothesis in part (d) require clarification.
\subsection{Sources}
\begin{itemize}
\item Supplied local UnsolvedMath registry record: permanent ID \(700722\), exact source code \texttt{catalogue entry}, title \emph{Kourovka Notebook Problem 11.29}. Attribution: N. D. Gupta; proposed in 1990.
\item \emph{The Kourovka Notebook: Unsolved Problems in Group Theory}, 11th issue (1990), Problem 11.29; the supplied record describes an update in the 21st issue (2026), identified as \texttt{arXiv:1401.0300v48}. The supplied collection description names E. I. Khukhro and V. D. Mazurov as editors. These source descriptions have not been independently checked.
\item Supplied Archive of Solved Problems attribution for the affirmative answer to part (f): N. D. Gupta and Yu. V. Kuz'min, \emph{J. Pure Appl. Algebra}, volume 78, no. 1 (1992), pp. 165--172. This citation is reproduced as supplied; no article title or external link is provided.
\item Supplied status review dated \texttt{2026-10-16}: canonical status \texttt{partially\_solved}, evidence label \texttt{published}; its note places solved parts in the archive and remaining parts in the main section. This is supplied dated curation, not an independently verified literature assessment.
\end{itemize}
\end{document}
